\documentclass[11pt]{article}
\usepackage[margin=25mm]{geometry}
\usepackage{fontspec,unicode-math}
\setmainfont{Latin Modern Roman}
\setmathfont{Latin Modern Math}
\usepackage{amsmath,amssymb,mathtools}
\usepackage{xurl,hyperref}
\hypersetup{colorlinks=true,urlcolor=blue,linkcolor=blue}
\setcounter{secnumdepth}{0}
\setlength{\parindent}{0pt}
\setlength{\parskip}{5pt}
\setlength{\emergencystretch}{3em}
\newcommand{\sref}[2]{\hyperlink{source:#1:#2}{[S#2]}}
\newcommand{\cataloguescope}[1]{\paragraph{Source question.}#1}
\newcommand{\floor}[1]{\left\lfloor#1\right\rfloor}
\begin{document}
\section{Geometry, residues, and scale thresholds for prescribed-prefix prime products}\label{statement}

\subsection{Definitions and mathematical statement}

\paragraph{Shared notation.}Standard number systems are denoted by $\mathbb N,\mathbb Z,\mathbb Q,\mathbb R,\mathbb C$; index ranges, computational models and specialized notation follow the cited problem source and the definitions below.


An integer with d decimal digits lies in $[10^{d-1},10^d)$. Saying n is the first k digits of a product means its decimal expansion begins with that k-digit string. The source fixes k=10000 and successively asks for prime factors of 10000, 7500, and 6000 digits.
\paragraph{Statement recorded in the supplied source.}
Question 11. Given a (random) number n of 10000 digits, it may be impossible to find 5000 digit primes p, q such that n = pq. Much easier: find some 10000 digit primes

p, q such that n is the first 10000 digits of pq. Instead, find 7500 digit primes p, q such that

n is the first 10000 digits of pq. (Coppersmith) Therefore, do this for 6000 digit primes. (Bernstein) \sref{statement}{2}
\paragraph{Scope qualification.}
Prescribed leading digits are the objective, rather than exact factorization of the given integer; the source does not specify a worst-case guarantee or success probability for the random-input wording.
\paragraph{Source status.}
The supplied catalogue labels this record partially solved. Its source-status metadata was checked against the supplied snapshot on 2026-10-11. This import does not independently establish current open status or mathematical correctness.
\subsection{Short English statement}
Can one construct primes of the requested lengths whose product begins with a prescribed ten-thousand-digit integer?
\subsection{Sources}
\begin{itemize}
\item[\hypertarget{source:statement:1}{S1}] Ulam AI and the UnsolvedMath Contributors, UnsolvedMath: A Curated Collection of Open Mathematics Problems, catalogue entry. Catalogue curation and metadata: CC BY 4.0. The supplied source question is retained; definitions and the English summary are editorial additions.
\url{https://huggingface.co/datasets/ulamai/UnsolvedMath}.
\item[\hypertarget{source:statement:2}{S2}] AIM Workshop Problem Lists, original problem source linked by the catalogue
\url{https://aimath.org/WWN/primesinp/primesinp.pdf}.
\end{itemize}
\paragraph{Preparation and review.}
Original question attributed to the cited source and UnsolvedMath contributors. Codex (GPT-6; exact serving revision not exposed) assisted with definitions, English summaries, source inspection and formatting. No new human mathematical authorship or independent proof verification is asserted. Catalogue identity checked against the complete supplied registry and existing repository statements; known flagged repetitions and obvious identical conjectures excluded. Full mathematical-equivalence/literature review has not been performed. Statement preparation and final review followed contribution prompts C and D; automated checks verify structure and identity, not open status.
\label{end:statement}
\end{document}
